% !TeX root = part6.tex
\documentclass[../final.tex]{subfiles}
\begin{document}
% Условия: ТиСФ2_семинар_6.pdf. Конспект: Desktop/1.jpg--3.jpg.
% Незавершённое вычисление скорости звука дополнено.
\practice{6}{23.10.2026}{Фотоны}

\rtheory{Равновесный фотонный газ}
Полагаем $k_B=1$: $T$ имеет размерность энергии. У фотона
$\varepsilon=cp=\hbar ck$ и две поперечные поляризации.
Число фотонов не сохраняется: стенки могут их поглощать и испускать.
Поэтому в тепловом равновесии $\mu_\gamma=0$, и
\[
 \overline n_\gamma(\varepsilon)=\frac1{e^{\varepsilon/T}-1},\qquad
 D_\gamma(\varepsilon)=\frac{V\varepsilon^2}{\pi^2\hbar^3c^3}.
\]
Используем интегралы
\begin{align*}
 \int_0^\infty\frac{x^{s-1}\,dx}{e^x-1}&=\Gamma(s)\zeta(s),\\
 \int_0^\infty\frac{x^{s-1}\,dx}{e^x+1}
 &=(1-2^{1-s})\Gamma(s)\zeta(s).
\end{align*}
Число фотонов и энергия равны
\[
 N_\gamma=\frac{2\zeta(3)}{\pi^2}V
 \left(\frac{T}{\hbar c}\right)^3,\qquad
 E_\gamma=aVT^4,\qquad a=\frac{\pi^2}{15\hbar^3c^3}.
\]
Большой термодинамический потенциал, давление и энтропия:
\[
 \Omega_\gamma=-\frac13aVT^4,\qquad
 P_\gamma=\frac13aT^4,\qquad S_\gamma=\frac43aVT^3.
\]
В обычных единицах плотность энергии равна $a(k_BT)^4$.
Коэффициент $a$ здесь относится к \emph{энергии}; коэффициент
перед $T^4$ в давлении в три раза меньше.

\clearpage
\task[1]{Электроны, позитроны и фотоны при высокой температуре}
\condition
При $T\gg m_ec^2$ найти отношения концентраций электронов,
позитронов и фотонов, а также отношения их энергий.
Концентрации электронов и позитронов считать одинаковыми.

\subsection*{Химические потенциалы}
Реакция $e^-+e^+\rightleftarrows2\gamma$ даёт
$\mu_-+\mu_+=2\mu_\gamma=0$.
Равенство концентраций при одинаковой дисперсии требует
$\mu_-=\mu_+$, следовательно $\mu_-=\mu_+=0$.
В ультрарелятивистском пределе для обоих сортов частиц
$\varepsilon\simeq cp$. Каждый имеет два спиновых состояния,
поэтому плотность состояний совпадает с фотонной:
\[
 D_-(\varepsilon)=D_+(\varepsilon)=D_\gamma(\varepsilon).
\]

\subsection*{Концентрации и энергии}
Для каждого сорта фермионов
\begin{align*}
 N_- = N_+&=\frac{VT^3}{\pi^2\hbar^3c^3}
 \int_0^\infty\frac{x^2\,dx}{e^x+1}
 =\frac34N_\gamma,\\
 E_- = E_+&=\frac{VT^4}{\pi^2\hbar^3c^3}
 \int_0^\infty\frac{x^3\,dx}{e^x+1}
 =\frac78E_\gamma.
\end{align*}
\answer
\[
 \boxed{n_-:n_+:n_\gamma=3:3:4},\qquad
 \boxed{E_-:E_+:E_\gamma=7:7:8}.
\]
Для электронов и позитронов вместе:
\[
 \frac{n_-+n_+}{n_\gamma}=\frac32,\qquad
 \frac{E_-+E_+}{E_\gamma}=\frac74,\qquad
 E_{\mathrm{tot}}=\frac{11}{4}E_\gamma.
\]
Числа $3/4$ и $7/8$ относятся к одному сорту фермионов,
а $3/2$ и $7/4$ --- к сумме двух сортов.

\clearpage
\task[2]{Влияние равновесного излучения на скорость звука}
\condition
Найти изменение скорости звука
$c_s^2=(\partial P/\partial\rho)_s$ в газе частиц массы $m$
и концентрации $n$ из-за фотонов при $T\ll mc^2$.

\subsection*{Уравнение состояния и условие адиабаты}
Считаем вещество классическим одноатомным идеальным газом,
а фотоны --- находящимися с ним в локальном тепловом равновесии
на времени звукового колебания. Плотность массы
$\rho=mn$. Инерцией излучения пренебрегаем, что требует
$aT^4\ll nmc^2$. Давление и полная энтропия:
\[
 P=nT+\frac13aT^4,\qquad
 S=N\left[\frac52-\ln\left(\frac{n\lambda_T^3}{g}\right)\right]
 +\frac43aVT^3,\qquad
 \lambda_T=\sqrt{\frac{2\pi\hbar^2}{mT}}.
\]
Здесь $g$ --- постоянное внутреннее вырождение атомов.
Производная в определении $c_s$ берётся при постоянной энтропии
на частицу $s=S/N$ (эквивалентно --- на единицу массы).
Обозначим
\[
 x=\frac{4aT^3}{3n}=4\frac{P_\gamma}{P_{\mathrm{gas}}}.
\]
Дифференцируя $s$, получаем
\[
 ds=\left(\frac32+3x\right)\frac{dT}{T}
 -(1+x)\frac{dn}{n}.
\]
Таким образом, на адиабате
\[
 \left(\frac{dT}{dn}\right)_s
 =\frac{T}{n}\frac{1+x}{3/2+3x}.
\]
При $x\to0$ это обычный закон $T\propto n^{2/3}$;
при $x\gg1$ --- радиационная адиабата $T\propto n^{1/3}$.

\clearpage
\subsection*{Производная давления}
Поскольку $P=P(n,T)$,
\begin{align*}
 c_s^2&=\frac1m\left[
 \left(\frac{\partial P}{\partial n}\right)_T
 +\left(\frac{\partial P}{\partial T}\right)_n
 \left(\frac{dT}{dn}\right)_s\right]\\
 &=\frac{T}{m}\left[1+\frac{(1+x)^2}{3/2+3x}\right].
\end{align*}
\answer
\[
 \boxed{c_s^2=\frac{T}{m}\left[\frac53+
 \frac{2x^2}{3(1+2x)}\right]},\qquad x=\frac{4aT^3}{3n}.
\]
Сравниваем с газом без фотонов при тех же $n,T$:
\[
 c_{s0}^2=\frac{5T}{3m},\qquad
 \boxed{\frac{c_s^2-c_{s0}^2}{c_{s0}^2}
 =\frac{2x^2}{5(1+2x)}}.
\]
Следовательно,
\[
 \frac{c_s-c_{s0}}{c_{s0}}
 =\sqrt{1+\frac{2x^2}{5(1+2x)}}-1
 \simeq\frac{x^2}{5}\quad(x\ll1).
\]
В первом порядке по малой доле излучения поправка сокращается:
фотоны одновременно увеличивают давление и теплоёмкость смеси.
При $x\gg1$ ведущий вклад равен $4aT^4/(9mn)$,
пока сохраняется условие малой инерции излучения.

Если температура задана в кельвинах, в этих формулах следует
заменить $T$ на $k_BT$, в том числе в определении $x$.
Для двухатомного газа или при нарушении теплового равновесия
теплоёмкость и адиабата будут другими.
\end{document}
